% ═══════════════════════════════════════════════════════════════
% AB-02b: tener + edad (Alter angeben)
% Abgeleitet aus LE-02b (Score 9.3/10)
% ═══════════════════════════════════════════════════════════════

\documentclass[11pt, a4paper]{article}

% ═══════════════════════════════════════════════════════════════
% SW-MODUS TOGGLE
% ═══════════════════════════════════════════════════════════════
% Setze \swmodustrue für Schwarz-Weiß-Druck
\newif\ifswmodus
\swmodusfalse    % Standard: Farbe

\usepackage[utf8]{inputenc}
\usepackage[T1]{fontenc}
\usepackage[ngerman]{babel}
\usepackage{helvet}
\renewcommand{\familydefault}{\sfdefault}

\usepackage[margin=2cm]{geometry}
\usepackage{enumitem}
\usepackage{tabularx}
\usepackage{booktabs}

\usepackage{xcolor}
\usepackage{tcolorbox}
\tcbuselibrary{skins}

\usepackage{fancyhdr}
\usepackage{lastpage}
\usepackage{needspace}

% Farben - Basis
\definecolor{spanisch-color}{HTML}{c41e3a}
\definecolor{grau-color}{HTML}{888888}
\definecolor{hellgrau-color}{HTML}{f5f5f5}
\definecolor{orange-color}{HTML}{b35c00}

% SW-Farben
\definecolor{sw-dark}{HTML}{000000}
\definecolor{sw-gray}{HTML}{666666}
\definecolor{sw-white}{HTML}{ffffff}

% Bedingte Farbzuweisung
\ifswmodus
    \colorlet{spanisch}{sw-dark}
    \colorlet{grau}{sw-gray}
    \colorlet{hellgrau}{sw-white}
    \colorlet{orange}{sw-dark}
\else
    \colorlet{spanisch}{spanisch-color}
    \colorlet{grau}{grau-color}
    \colorlet{hellgrau}{hellgrau-color}
    \colorlet{orange}{orange-color}
\fi

\newcommand{\fachfarbe}{spanisch}

% Variablen
\newcommand{\thema}{tener + edad (Alter angeben)}
\newcommand{\sequenz}{02}
\newcommand{\block}{b}
\newcommand{\fach}{Spanisch}
\newcommand{\klasse}{7}
\newcommand{\maxpunkte}{20}

% Kopf-/Fußzeile
\pagestyle{fancy}
\fancyhf{}
\renewcommand{\headrulewidth}{0.4pt}
\renewcommand{\footrulewidth}{0.4pt}
\lhead{\fach{} -- Klasse \klasse}
\chead{AB-\sequenz\block}
\rhead{}
\lfoot{}
\rfoot{Seite \thepage\ von \pageref{LastPage}}

% Befehle
\newcommand{\punktebox}[1]{\hfill\fbox{\small \_/#1}}
\newcommand{\luecke}[1]{\rule{#1}{0.4pt}}
\newcommand{\es}[1]{\textcolor{\fachfarbe}{\textbf{#1}}}

% Merkkasten (SW-kompatibel: schwebender Titel)
\newtcolorbox{merkkasten}[1][]{
    colback=hellgrau,
    colframe=\fachfarbe,
    colbacktitle=white,
    coltitle=\fachfarbe,
    boxrule=1pt,
    arc=0pt,
    left=8pt, right=8pt, top=8pt, bottom=8pt,
    fonttitle=\bfseries,
    attach boxed title to top left={yshift=-2mm, xshift=5mm},
    boxed title style={boxrule=0pt, colback=white},
    title=#1
}

% Hinweis (SW-kompatibel)
\newtcolorbox{hinweis}{
    colback=white,
    colframe=orange,
    colbacktitle=white,
    coltitle=orange,
    boxrule=1pt,
    arc=0pt,
    left=5pt, right=5pt, top=5pt, bottom=5pt,
    fonttitle=\bfseries,
    attach boxed title to top left={yshift=-2mm, xshift=5mm},
    boxed title style={boxrule=0pt, colback=white},
    title=Hinweis
}

% Aufgabe (kein Seitenumbruch innerhalb)
\newenvironment{aufgabe}[1]{%
    \needspace{5\baselineskip}%
    \nopagebreak[4]%
    \vspace{10pt}
    \noindent\textbf{\textcolor{\fachfarbe}{Aufgabe #1}}
    \nopagebreak[4]%
    \vspace{5pt}
    \nopagebreak[4]%
    \begin{list}{}{\leftmargin=0pt}
    \item[]
}{%
    \end{list}
}

% ═══════════════════════════════════════════════════════════════
\begin{document}

% Kopfbereich
\begin{center}
    {\Large\bfseries\textcolor{\fachfarbe}{Arbeitsblatt: \thema}}
\end{center}

\vspace{5pt}

\noindent
\begin{tabularx}{\textwidth}{|X|X|X|}
\hline
\textbf{Name:} \luecke{4cm} &
\textbf{Klasse:} \klasse &
\textbf{Datum:} \luecke{2cm} \\
\hline
\end{tabularx}

\vspace{15pt}

% ═══════════════════════════════════════════════════════════════
% MERKKASTEN 1: Verb tener ($\rightarrow$ FZ1)
% ═══════════════════════════════════════════════════════════════

\begin{merkkasten}[Merkkasten: Das Verb \textit{tener} (haben)]

\begin{center}
\begin{tabular}{|l|l|l|}
\hline
\textbf{Person} & \textbf{Pronomen} & \textbf{tener} \\
\hline
1. Sg. & yo (ich) & \luecke{2.5cm} \\
\hline
2. Sg. & tú (du) & \luecke{2.5cm} \\
\hline
3. Sg. & él/ella (er/sie) & \luecke{2.5cm} \\
\hline
1. Pl. & nosotros/-as (wir) & \luecke{2.5cm} \\
\hline
\end{tabular}
\end{center}

\vspace{0.5em}

\textbf{Alter angeben:} \es{¿Cuántos años tienes?} = Wie alt bist du? (wörtlich: Wie viele Jahre hast du?)

\vspace{0.3em}

$\rightarrow$ \es{Tengo ... años.} = Ich bin ... Jahre alt. (wörtlich: Ich habe ... Jahre.)

\vspace{0.5em}
\textit{Übertrage die Formen aus dem Lernpfad oder von der Tafel!}
\end{merkkasten}

\vspace{10pt}

% ═══════════════════════════════════════════════════════════════
% MERKKASTEN 2: Zahlen 1-20 ($\rightarrow$ FZ2)
% ═══════════════════════════════════════════════════════════════

\begin{merkkasten}[Merkkasten: Zahlen 1-20 (Los números)]

\begin{center}
\begin{tabular}{|c|l|c|l|c|l|c|l|}
\hline
1 & uno & 6 & seis & 11 & once & 16 & dieciséis \\
\hline
2 & dos & 7 & siete & 12 & doce & 17 & diecisiete \\
\hline
3 & tres & 8 & ocho & 13 & trece & 18 & dieciocho \\
\hline
4 & cuatro & 9 & nueve & 14 & catorce & 19 & diecinueve \\
\hline
5 & cinco & 10 & diez & 15 & quince & 20 & veinte \\
\hline
\end{tabular}
\end{center}

\vspace{0.3em}
\textit{Achtung: 11-15 sind besondere Formen! 16-19 = dieci + Einer}
\end{merkkasten}

\vspace{15pt}

% ═══════════════════════════════════════════════════════════════
% AUFGABE 1: tener einsetzen ($\rightarrow$ FZ1, AFB I)
% ═══════════════════════════════════════════════════════════════

\begin{aufgabe}{1}
Setze die richtige Form von \textit{tener} ein. \punktebox{4}
\end{aufgabe}

\begin{enumerate}[label=\alph*)]
    \item Yo \luecke{2.5cm} 12 años.
    \item Mi hermana \luecke{2.5cm} 15 años.
    \item ¿Cuántos años \luecke{2.5cm} tú?
    \item Nosotros \luecke{2.5cm} un perro.
\end{enumerate}

\vspace{15pt}

% ═══════════════════════════════════════════════════════════════
% AUFGABE 2: Zahlen zuordnen ($\rightarrow$ FZ2, AFB I)
% ═══════════════════════════════════════════════════════════════

\begin{aufgabe}{2}
Verbinde die spanischen Zahlen mit der richtigen Ziffer. \punktebox{4}
\end{aufgabe}

\begin{center}
\begin{tabular}{rcl}
\es{doce} & \luecke{3cm} & 8 \\[0.8em]
\es{diecisiete} & \luecke{3cm} & 15 \\[0.8em]
\es{ocho} & \luecke{3cm} & 12 \\[0.8em]
\es{quince} & \luecke{3cm} & 17 \\
\end{tabular}
\end{center}

\vspace{15pt}

% ═══════════════════════════════════════════════════════════════
% AUFGABE 3: Alter angeben ($\rightarrow$ FZ3, AFB II)
% ═══════════════════════════════════════════════════════════════

\begin{aufgabe}{3}
Beantworte die Fragen auf Spanisch. \punktebox{6}
\end{aufgabe}

\begin{hinweis}
\textbf{Struktur:} Tengo/Tiene + [Zahl] + años.
\end{hinweis}

\vspace{0.5em}

\noindent
a) \es{¿Cuántos años tienes?}

\vspace{0.3em}
\noindent\rule{\textwidth}{0.4pt}\vspace{8pt}

\noindent
b) \es{¿Cuántos años tiene tu madre?}

\vspace{0.3em}
\noindent\rule{\textwidth}{0.4pt}\vspace{8pt}

\noindent
c) \es{¿Cuántos años tiene tu hermano/hermana?} (Falls du keine Geschwister hast: \textit{un amigo/una amiga})

\vspace{0.3em}
\noindent\rule{\textwidth}{0.4pt}

\vspace{15pt}

% ═══════════════════════════════════════════════════════════════
% AUFGABE 4: Familie beschreiben ($\rightarrow$ FZ4/FZ5, AFB II/III)
% ═══════════════════════════════════════════════════════════════

\begin{aufgabe}{4}
Beschreibe deine Familie. Schreibe mindestens 3 Sätze mit Namen und Alter. \punktebox{6}
\end{aufgabe}

\begin{hinweis}
\textbf{Beispiel:} Mi padre se llama Carlos. Tiene 45 años.
\end{hinweis}

\noindent\rule{\textwidth}{0.4pt}\vspace{8pt}
\noindent\rule{\textwidth}{0.4pt}\vspace{8pt}
\noindent\rule{\textwidth}{0.4pt}\vspace{8pt}
\noindent\rule{\textwidth}{0.4pt}\vspace{8pt}
\noindent\rule{\textwidth}{0.4pt}\vspace{8pt}
\noindent\rule{\textwidth}{0.4pt}

\vspace{20pt}
\hrule
\vspace{10pt}

\begin{flushright}
\textbf{Punkte:} \luecke{1cm} / \maxpunkte
\end{flushright}

\end{document}
